\documentclass[12pt, hyperref={bookmarks=false}, show notes]{beamer}
% Switch to this for handouts
%\documentclass[12pt,handout]{beamer}

% Text
	\usepackage[T1]{fontenc}
  \usepackage{emerald}
	\usepackage[utf8]{inputenc}
	\usepackage[english]{babel}
	% \usepackage[bitstream-charter]{mathdesign} % Serif font (Charter BT).
	\usepackage[scaled=0.84]{DejaVuSansMono} % Monospaced font.
	\def\sfdefault{SourceSansPro-TLF} % Sans serif font.
	\usepackage{textcomp}

% Maths
  \usepackage{amsmath}
  \usepackage{mathtools}
  \usepackage{siunitx}

% Graphics
  \usepackage{graphicx}
  \usepackage[caption=false]{subfig}
	\usepackage{tikz}
  \usepackage{pgfplots}
  \pgfplotsset{compat=1.10}
	% ADD TIKZ LIBRARIES
  \usetikzlibrary{calc}
  \usetikzlibrary{arrows.meta}
  %\usepackage{tikz-qtree}
  \usetikzlibrary{decorations.pathmorphing}
  \usetikzlibrary{matrix,shapes,positioning}
  \usepgfplotslibrary{external}
  \tikzexternalize[prefix=gfx/tikz/]
  \tikzexternaldisable % Disable by default
  \usepackage{pgfgantt}
  \usepackage{mdframed}
  \newmdenv [align=center, backgroundcolor=color3!10, linecolor=color3!50, linewidth=1pt]{mtjbox}

	\usepackage{xcolor}
    \definecolor{color1}{cmyk}{100,50,0,0}   % blue
    \definecolor{color2}{cmyk}{0,80,100,0}   % vermillion
    \definecolor{color3}{cmyk}{97,0,75,0}    % blueish green
    \definecolor{color4}{RGB}{204,121,167}    % reddish purple
    \definecolor{color5}{RGB}{230,159,0}   % orange
    \usepackage{colortbl}
% Misc
\usepackage{booktabs}
\usepackage{enumerate}
\usepackage{pdfpages}
\usepackage{pgfpages}
%\usepackage{setspace}
%\usepackage{multimedia}

% Box colors.
\definecolor{clframe}{gray}{0.75}
\definecolor{clshade}{gray}{0.95}
\definecolor{clcodeshade}{gray}{0.93}

% Code colors.
\definecolor{javared}{rgb}{0.6,0,0}            % for strings
\definecolor{javagreen}{rgb}{0.25,0.5,0.35}    % comments
\definecolor{javapurple}{rgb}{0.6,0,0.5}       % keywords
\definecolor{javadocblue}{rgb}{0.25,0.35,0.75} % javadoc

%More colors
\definecolor{colorA}{HTML}{5DA5DA} % blue
\definecolor{colorB}{HTML}{FAA43A} % orange
\definecolor{colorC}{HTML}{60BD68} % green
\definecolor{colorD}{HTML}{F17CB0} % pink
\definecolor{colorE}{HTML}{B2912F} % brown
\definecolor{colorF}{HTML}{B276B2} % purple
\definecolor{colorG}{HTML}{DECF3F} % yellow
\definecolor{colorH}{HTML}{F15854} % red
\definecolor{colorI}{HTML}{4D4D4D} % gray

\colorlet{colorAshade}{colorA!80!black}
\colorlet{colorBshade}{colorB!80!black}
\colorlet{colorCshade}{colorC!80!black}
\colorlet{colorDshade}{colorD!80!black}
\colorlet{colorEshade}{colorE!80!black}
\colorlet{colorFshade}{colorF!80!black}
\colorlet{colorGshade}{colorG!80!black}
\colorlet{colorHshade}{colorH!80!black}
\colorlet{colorIshade}{colorI!80!black}

\usepackage{listings}

\lstdefinestyle{javastyle}{%
  language=Java,
  keywordstyle=\color{javapurple}\bfseries,
  stringstyle=\color{javared},
  commentstyle=\color{javagreen},
  morecomment=[s][\color{javadocblue}]{/**}{*/}
}

\lstnewenvironment{javacode}[1][]{%
  \lstset{style=javastyle, #1}
}{}

\lstset{
  columns=flexible,
  escapechar=",
  escapeinside={(*@}{@*)},
  basicstyle=\ttfamily\footnotesize,
  frame=single, backgroundcolor=\color{clshade}, rulecolor=\color{clframe},
  framerule=\fboxrule, xleftmargin=3.4pt, xrightmargin=3.4pt, belowskip=\smallskipamount,
  captionpos=b,
  numbers=left,
  numberstyle=\scriptsize,
  numbersep=7pt,
  tabsize=2,
  keepspaces=true,
  showspaces=false,
  showstringspaces=false,
  breaklines=true,
  breakatwhitespace=true,
  postbreak=\raisebox{0ex}[0ex][0ex]{\ensuremath{\color{red}\hookrightarrow\space}}
}



\setbeamertemplate{navigation symbols}{}
%\setbeamertemplate{bibliography item}[text]

%\usepackage{caption}
\usetheme{/amsterdam}
\date{June 15, 2015}

%Uncomment for handouts: (remember to uncomment in line 1+2, too)
%\usepackage{beamertheme/handoutWithNotes}
%\pgfpagesuselayout{4 on 1 with notes}[a4paper,border shrink=5mm]

% Comment for handouts or slides without notes:
\setbeameroption{show notes on second screen=right}
\setbeamertemplate{note page}{%
  \insertnote%
}

% Table of content dybde (0-index)
%\setcounter{tocdepth}{1}

% BibLaTeX
%\usepackage{csquotes}
%\usepackage[
%backend=bibtex,
%citestyle=numeric,
%bibstyle=numeric,
%maxcitenames=3,
%maxbibnames=99,
%url=true]{biblatex}
%\addbibresource{../rapport/references/refs.bib}
%\addbibresource{extrasources.bib}
%\usepackage{../style/biblatex_custom_formatting}

\graphicspath{{gfx/}{../report/graphics/}{../report/graphics/tikz}}

\begin{document}

%\captionsetup[figure]{font=small,singlelinecheck=off,justification=raggedright}

\title[Improving Wikipedia Linking: A Feature Learning Approach]{Improving Wikipedia Linking: A Feature Learning Approach}
% \author[\insertframenumber /\inserttotalframenumber]{sw609f15}
\makeatletter
  \defbeamertemplate*{footline}{myminiframes theme}
  {%
    \begin{beamercolorbox}[colsep=1.5pt]{upper separation line foot}
    \end{beamercolorbox}
    \hbox{%
      \begin{beamercolorbox}[wd=.75\paperwidth,ht=2.5ex,dp=1.125ex,%
        leftskip=.3cm,rightskip=.3cm plus1fil]{title in head/foot}%
        \leavevmode{\usebeamerfont{title in head/foot}\insertshorttitle}%
      \end{beamercolorbox}%
      \begin{beamercolorbox}[wd=.15\paperwidth,ht=2.5ex,dp=1.125ex,%
        leftskip=.3cm,rightskip=.3cm]{title in head/foot}%
        {\usebeamerfont{author in head/foot}\usebeamercolor[fg]{author in head/foot}\insertshortauthor}
      \end{beamercolorbox}%
      \begin{beamercolorbox}[wd=.1\paperwidth,ht=2.5ex,dp=1.125ex,%
        leftskip=.3cm,rightskip=.3cm,center]{title in head/foot}%
        {\usebeamerfont{author in head/foot}\usebeamercolor[fg]{author in head/foot}\insertframenumber/\inserttotalframenumber}%
      \end{beamercolorbox}%
    }%
    \begin{beamercolorbox}[colsep=1.5pt]{lower separation line foot}
    \end{beamercolorbox}
  }
\makeatother

% Adding front slide.
{
\setbeamercolor{background canvas}{bg=}
\includepdf[pages=1]{frontslide.pdf}
}
\note{
    Noter, noter, noter!
}

\begin{frame}
  \frametitle{Indhold}
  \tableofcontents
  \note{
  \begin{itemize}
    \item Notes
    \item notes
    \item notes..
  \end{itemize}
}
\end{frame}


\chapter{Introduction}
The World Wide Web is built upon the idea of \emph{creating a space in which anything could be linked to anything} \todo{insert cite: Weaving the Web, chapter 1 by Tim Berners-Lee}. Thus, hyperlinking can be seen as a vital component of the World Wide Web, and even a critical component to its success. Hyperlinking allows websites to be easily navigationable, enabling users to quickly find relevant or in-depth information.

A wide array of research that tries to find links.

\todo{motivation and users}

\section{Problem Statement}
\dummy

\section{Project Goals}
\todo{Success criteria (to be evaluated in the conclusion)}

\section{Compliance with Study Regulation}
\todo{Maybe a description of how the project fits within the study regulation}

\section{Report Organization}
This report describes and documents the design, implementation, evaluation and development method of the system described in this chapter. Immediately following this introduction, \chapterref{chap:analysis} ... \dummy ... \chapterref{chap:devmethodreflection} contains a discussion and reflection of the development process used throughout this project.
\section[Theory]{Theory}
% Første parameter i [] er tekst i header. {} er i indholdsfortegnelsen.

% Slide med emneoverskrift.
\begin{frame}
  \frametitle{}
  \begin{center}
    {\Huge node2vec}
  \end{center}
\end{frame}
\note{
  \begin{itemize}
		\item Notes...
  \end{itemize}
}

% Normal slide:
\begin{frame}
    \frametitle{Some Example Title}
    \framesubtitle{Some example subtitle}
    \centering
    Some text, content, etc.
\end{frame}
\note{
	\begin{itemize}
    \item Notes here...
	\end{itemize}
}
\section[Wronglywent]{What Went Wrong?}
% Første parameter i [] er tekst i header. {} er i indholdsfortegnelsen.

% Slide med emneoverskrift.
\begin{frame}
  \frametitle{}
  \begin{center}
    {\Huge What Went Wrong?}
  \end{center}
\end{frame}
\note{
  \begin{itemize}
		\item Notes...
  \end{itemize}
}

% Normal slide:
\begin{frame}
    \frametitle{Some Example Title}
    \framesubtitle{Some example subtitle}
    \centering
    Some text, content, etc.
\end{frame}
\note{
	\begin{itemize}
    \item Notes here...
	\end{itemize}
}
\section[New Work]{New Work}
% Første parameter i [] er tekst i header. {} er i indholdsfortegnelsen.

% Slide med emneoverskrift.
\begin{frame}
  \frametitle{}
  \begin{center}
    {\Huge New Work}
  \end{center}
\end{frame}
\note{
  \begin{itemize}
		\item Notes...
  \end{itemize}
}

% Normal slide:
\begin{frame}
    \frametitle{Some Example Title}
    \framesubtitle{Some example subtitle}
    \centering
    Some text, content, etc.
\end{frame}
\note{
	\begin{itemize}
    \item Notes here...
	\end{itemize}
}
\section[Future Work]{Future Work}
% Første parameter i [] er tekst i header. {} er i indholdsfortegnelsen.

% Slide med emneoverskrift.
\begin{frame}
  \frametitle{}
  \begin{center}
    {\Huge Future Work}
  \end{center}
\end{frame}
\note{
  \begin{itemize}
		\item Notes...
  \end{itemize}
}

% Normal slide:
\begin{frame}
    \frametitle{Future Work}
    \framesubtitle{Reinforcement Learning}
    \centering
    Some text, content, etc.
\end{frame}
\note{
	\begin{itemize}
    \item Notes here...
	\end{itemize}
}

\section*{}

\begin{frame}
  \frametitle{Overview}
  \tableofcontents
\end{frame}

\end{document}
